\section{Conclusiones: resolución afirmativa de la hipótesis del continuo}
\label{sec:conclusiones-ch}
\begingroup
\fontsize{9.7}{11.8}\selectfont
\setlength{\parskip}{1.5pt}
\setlength{\abovedisplayskip}{4pt}
\setlength{\belowdisplayskip}{4pt}

La composición
\[
 \mathrm{APP}\longrightarrow\mathrm{TRIT}\longrightarrow\mathrm{TPK}
 \longrightarrow X_\infty^{\mathrm{enr}}
 \longrightarrow\mathcal C_{\mathrm{cont}}^{\mathrm{disc}}
\]
genera un único sistema compatible de historias dotadas de valor, hoja,
orientación, residuo, cociente, acarreo, frontera, ruta, memoria y
supervivencia. Su proyección arquimediana produce la recta
\(\mathbb R_{\mathrm{HMT}}^{\mathfrak G}\); su proyección incidencial conserva
la información de frontera que la publicación del valor no retiene. La recta
real es, por tanto, una salida de la construcción y no un dominio previamente
introducido para seleccionar sus estados.

El resultado cardinal principal es
\begin{equation}
 \boxed{
 \mathfrak G_{\mathrm{HMT}}(h,m_\infty)
 \models 2^{\aleph_0}=\aleph_1 .}
 \label{eq:conclusion-decision-continuo}
\end{equation}
La igualdad no se presupone en ninguna etapa generativa. La cobertura
exhaustiva de cilindros establece la equipotencia
\[
 \mathbb R_{\mathrm{HMT}}^{\mathfrak G}
 \approx\mathcal P^{\mathfrak G}(\omega).
\]
La condensación genealógica asigna a cada real interno un ordinal estrictamente
menor que \(\omega_1^{\mathfrak G}\) y construye
\[
 j_{\mathrm{HMT}}:
 \mathbb R_{\mathrm{HMT}}^{\mathfrak G}\hookrightarrow
 \omega_1^{\mathfrak G}.
\]
La codificación de los buenos órdenes numerables y su levantamiento ternario
construyen la inyección en sentido contrario
\[
 k_{\mathrm{HMT}}:
 \omega_1^{\mathfrak G}\hookrightarrow
 \mathbb R_{\mathrm{HMT}}^{\mathfrak G}.
\]
El teorema de Cantor--Bernstein concluye la igualdad cardinal de
\eqref{eq:conclusion-decision-continuo}; no la suministra como entrada.

El cuantificador del resultado alcanza todos los subconjuntos de la recta que
pertenecen a la clausura genealógica completa. Para cada
\(A\in\mathcal P^{\mathfrak G}
(\mathbb R_{\mathrm{HMT}}^{\mathfrak G})\), la cofinalidad demostrada impone
la dicotomía
\[
 |A|\leq\aleph_0
 \qquad\text{o}\qquad
 |A|=|\mathbb R_{\mathrm{HMT}}^{\mathfrak G}|.
\]
No existe, en el continuo íntegro generado por HMT, un cardinal estrictamente
intermedio entre \(\aleph_0\) y el cardinal de la recta. En consecuencia,
Holografía Modular Triádica resuelve afirmativamente la hipótesis del
continuo en su dominio completo de generación.

La finitud de cada nivel no restringe este alcance. La fibra
\(\mathbb F_3^6\), formada por \(3^6=729\) subcilindros de refinamiento,
describe la ramificación local; el cuantificador coinductivo actúa para toda
profundidad prescrita. Las ejecuciones de mil, diez mil o un millón de cifras
son testigos finitos de una ley parametrizada por la profundidad, no cotas de
la construcción. El retorno nonádico recupera la fase mientras el acarreo y la
memoria avanzan, de modo que el cierre observable no reinicia el estado
enriquecido.

La misma genealogía produce las publicaciones correlacionadas de la doble
realización. Entre ellas se encuentran la holonomía nonádica
\(\Gamma_9\), su lectura posterior de memoria reducida
\(\mathcal K_{\mathrm{ph}}\), la constante holográfica de acoplamiento
aritmético-geométrico \(K\), la cuatrirrelación
\((\pi,\varphi,e,\alpha)\) y el corredor
Paley--Witt--Golay--Leech--\(V^\natural\)--Monstruo. Son codominios y
lectores distintos de un mismo estado; ninguno se introduce desde valores
convencionales para forzar la comparación cardinal.

En particular, \(K\) es la representación integral canónica del registro
dodecafásico orientado. Su producción sigue la cadena
\[
 X_{\rm term}\longrightarrow\mathcal R_{12}(X_{\rm term})
 \longrightarrow(b^{90},b^{120},Q_{\rm TPK})
 \longrightarrow U_{\rm sgn},
 \qquad
 U_{\rm sgn}\xrightarrow{\;\mathcal H_{12}/4\;}K,
 \qquad
 K\xrightarrow{\;\mathcal H_{12}\;}U_{\rm sgn}.
\]
La equivalencia de Hadamard expresa una coordenatización reversible del
registro generado; no convierte a \(U_{\rm sgn}\), a \(K\) ni a su lectura
racional en datos externos. La lectura posicional de \(K\) interviene
posteriormente en la clausura de \(\alpha\); sus lecturas incidencial y
geométrica conservan soporte, dirección y reducción dimensional. Esta función
no debe confundirse con \(\mathcal K_{\mathrm{ph}}\), que pertenece a otro
tipo y se publica después de la holonomía.

La identificación
\[
 \mathfrak G_{\mathrm{HMT}}(h,m_\infty)
 =L[u_{h,m}]=L[h,m_\infty]
\]
es el reconocimiento final de una jerarquía ya construida. Hace explícita su
representación por constructibilidad relativa, pero no genera la recta, no
selecciona las inyecciones y no interviene en la reducción de cilindros. El
orden causal permanece APP--TRIT--TPK, estado enriquecido, estructura discreta
conjunta del continuo, publicación arquimediana e identificación posterior.

Quedan así demostradas conjuntamente la construcción del continuo HMT, la
equipotencia de su recta con la potencia de \(\omega\), las dos inyecciones
cardinales, la ausencia de cardinales intermedios en la clausura completa y la
igualdad \(2^{\aleph_0}=\aleph_1\). Ésta es la resolución afirmativa de la
hipótesis del continuo establecida por el teorema
\ref{thm:decision-continuo-hmt}.
\endgroup
